\section*{Chapter \ref{ch:ll:resilience} --- Resilience}

\begin{solution}{exo:ll:resilience:1}
The journal refuses the append: epoch 2 is older than 3. The entry is not recorded, so no replica applies it, and the orders
the old primary computed from it are never sent, since only journaled outputs are sent (and the gateway checks the epoch too).
\end{solution}

\begin{solution}{exo:ll:resilience:2}
The last \omterm{def:ll:protocols-i-fix:session}{heartbeat} was sent at some tick $h$ before the failure. Death is declared at the first tick at or after $h + 3$~ms: if
the failure came just after $h$, detection takes almost \qty{3}{\milli\second}; if it came just before the next tick, a little
over \qty{2}{\milli\second}. Hence the build's recovery times of 2.8 to \qty{3.1}{\milli\second}, reconnection included.
\end{solution}

\begin{solution}{exo:ll:resilience:3}
$46 \times 16 + 0 = 736$. Every replica must derive it from the journal so that reports, which name the order by this
identifier, are applied to the same order everywhere, and so that a resend under it is recognised by the venue as a copy.
\end{solution}

\begin{solution}{exo:ll:resilience:4}
$500 \times \qty{40}{\micro\second} \times 0.9 = 1.8\%$.
\end{solution}

\begin{solution}{exo:ll:resilience:5}
$37 \times 10^6 \times \qty{37}{\nano\second} \approx \qty{1.4}{\second}$ for the day, about \qty{0.2}{\second} for its last
hour. A hot backup needs no catch-up at all, and it has been exercised on every entry: a cold one is replaying code paths at
the worst moment.
\end{solution}

\begin{solution}{exo:ll:resilience:6}
A rejection with reason \texttt{D} (duplicate identifier), journaled like any report. The replica keeps the order open (the
venue has the original) and waits for the original's reports, which the session's replay delivers; it must not treat the
rejection as the end of the order.
\end{solution}

\begin{solution}{exo:ll:resilience:7}
With reports delayed, orders are still unreported when the backup takes over, even after its replay: ``fresh identifiers,
reconciled'' now resends them too, and duplicates those that the venue has already executed, at the stages ``sent'' and ``at
the venue''. Idempotency keys still produce none, because the venue recognises each copy by its identifier whatever the
backup knows.
\end{solution}

\begin{solution}{exo:ll:resilience:8}
Restarting from the configuration forgets the state: the position, the open orders at the venue, the orders in flight. The new
quotes are sent next to the old ones, which may still be live (unless cancel on disconnect removed them), and the position
limits are checked against a position of zero. The backup must rebuild the state from the journal, reconnect, reconcile, and
only then trade; and if the primary was only slow, nothing stops it from sending too: there is no fencing.
\end{solution}

\begin{solution}{pb:ll:resilience:1}
\begin{enumerate}
\item Journaled but not yet sent; sent and in flight; at the venue but not yet reported; reported. The new primary cannot tell,
in the second and third, whether the order left.
\item \qty{20}{\micro\second} on the wire, then \qty{20}{\micro\second} until the report is journaled: \qty{40}{\micro\second}.
\item A duplicate counts only if both copies executed; the other 19 orders found no liquidity at their limit and were
cancelled (immediate or cancel), so their copies were cancelled too.
\item The primary had not sent it, so the backup's resend is the only copy.
\item $46 / \qty{60}{\milli\second} \approx 767$ orders a second.
\item The dangerous windows add up to $27 \times \qty{40}{\micro\second} = \qty{1.08}{\milli\second}$ in
\qty{60}{\milli\second}: about 1.8\%.
\item $5\,000 \times \qty{40}{\micro\second} \times 27/46 \approx 11.7\%$.
\item The duplicate adds a second execution to a position the replica had already brought to its limit of 1\,000.
\item The venue reports every order within \qty{40}{\micro\second} and replays its reports on login, so after the replay
nothing sent is unaccounted for.
\item A venue that reports slowly (a resting order acknowledged late, a busy matching engine), or a session that cannot be
replayed (a new session, a venue without sequenced replay, the situation of the Tokyo participants in 2020).
\item The venue recognises the copy by its identifier; it does not matter what the new primary knows or in what order it acts.
\item Identifiers fixed when the order is created, derived from something every replica agrees on (the journal), and a venue
or gateway that rejects a reused identifier within the session.
\item \textbf{Named result.} Without idempotency keys and without reconciliation, a \omterm{def:ll:resilience:failover}{failover} duplicated an order at 54 of the
184 cut points (every executed order, whenever the failure fell between its departure and its report), about 1.8\% of failures
at a random time in the scenario, pushing the position to 1\,100 against a limit of 1\,000; with idempotency keys, none; and
recovery took 2.8 to \qty{3.1}{\milli\second}, almost all of it the \omterm{def:ll:protocols-i-fix:session}{heartbeat} timeout.
\item The \omterm{def:ll:protocols-i-fix:session}{heartbeat} timeout; shortening it declares slow primaries dead more often, which makes fencing indispensable.
\item It stops a primary that is not dead from recording and sending anything after the takeover; idempotency keys stop copies
of the same order, not new orders from a second primary.
\item The \omterm{def:ll:resilience:failover}{failover} (the switch to the second device) failed for an untested kind of failure; after the manual switch, the state
of the orders was not agreed between the exchange and its participants, and there was no tested procedure to reconcile and
restart.
\item So that every replica agrees on which identifier each order now has; otherwise a later \omterm{def:ll:resilience:failover}{failover} would not know about the
resend.
\item State equal to a fresh replay of the journal, position equal to the venue's, no execution twice with keys. Properties hold
whatever the failure point, so one check covers all 184 runs, and they say what correct means rather than what one run did.
\item Kill the backup's state first, restart it from the journal at every cut point, and check the same properties plus the
catch-up time.
\item One journal that every replica follows, identifiers derived from it, fencing of the old primary, and reconciliation before
trading resumes.
\end{enumerate}
\end{solution}

\begin{solution}{iq:ll:resilience:1}
A single component that numbers every input and journals it before anyone acts on it. It gives the system one order of events,
so that components are deterministic functions of the journal: they can be replicated, replayed and restarted.
\iqlookfor{total order, journal, determinism.}
\end{solution}

\begin{solution}{iq:ll:resilience:2}
It is hot (has applied the journal), so it fences the old primary with a new epoch, reconnects the venue session asking for the
missed reports, applies them, and resends the orders no report acknowledges under their original, journal-derived identifiers,
so that the venue rejects any copy.
\iqlookfor{fencing, replay of reports, idempotent identifiers.}
\end{solution}

\begin{solution}{iq:ll:resilience:3}
Two replicas acting as primary at once, usually after a takeover from a primary that was slow rather than dead. Prevent it with
fencing (an epoch that only grows, checked by the journal and the gateway) and, for electing the primary, a quorum.
\iqlookfor{epochs and quorum.}
\end{solution}

\begin{solution}{iq:ll:resilience:4}
Stop the strategy's new orders; log in asking for the reports missed; apply them; account for cancel on disconnect; reconcile
the open orders and position against the \omterm{def:ll:order-gateway-and-order-management:dropcopy}{drop copy}; resend or cancel what is unaccounted for with idempotent identifiers; then
resume.
\iqlookfor{replay, reconcile, then trade.}
\end{solution}

\begin{solution}{iq:ll:resilience:5}
Make failure a routine test: kill the primary at every cut point of scenarios in continuous integration and check properties
(state equals replay, positions equal the venue's, no duplicate); fail over in production on a schedule.
\iqlookfor{exhaustive cut points, properties, scheduled \omterm{def:ll:resilience:failover}{failovers}.}
\end{solution}

\begin{solution}{iq:ll:resilience:6}
A replicated sequencer and journal; hot replicas of the engine, the \omterm{def:ll:pre-trade-risk-and-kill-switches:gate}{risk gate} and the gateway as deterministic functions of it;
primary election with epochs; \omterm{def:ll:protocols-i-fix:session}{heartbeats} with a timeout of a few milliseconds; sessions with sequenced replay and idempotent
identifiers; reconciliation against drop copies. Tested by killing each machine at every cut point of recorded days.
\iqlookfor{journal, replicas, fencing, idempotency, tests.}
\end{solution}
